\begin{solapartat}
\punts{0,6}
\figura[0.7]{sol_fig1.png}
Si els nodes no estan ben identificats es restarà 0,2~p.\\
Si els ventres no estan ben identificats es restarà 0,2~p.

\punts{0,1} $L = \dfrac{\lambda_A}{2} \Rightarrow \lambda_A = 2L = 0,800\ \mathrm{m}$

\punts{0,1} $L = 2\,\dfrac{\lambda_B}{2} \Rightarrow \lambda_B = L = 0,400\ \mathrm{m}$

\punts{0,1} $L = 3\,\dfrac{\lambda_C}{2} \Rightarrow \lambda_C = \dfrac{2}{3}L = 0,267\ \mathrm{m}$

\punts{0,1} $L = 4\,\dfrac{\lambda_D}{2} \Rightarrow \lambda_D = \dfrac{1}{2}L = 0,200\ \mathrm{m}$

\punts{0,25} $v = \dfrac{\lambda_A}{T} = \lambda_A\cdot f_A = 0,8\cdot 120 = 96\ \mathrm{m/s}$

Aquest càlcul es pot fer amb qualsevol dels harmònics:
\[ v = \lambda_B\cdot f_B = \lambda_C\cdot f_C = \lambda_D\cdot f_D = 96\ \mathrm{m/s} \]
\end{solapartat}

\begin{solapartat}
% DUBTE: la pauta parla de «la longitud de la cavitat» tot i que l'enunciat tracta d'una corda; es transcriu tal com és.
\punts{0,65} Si la longitud de la cavitat disminueix a la meitat, llavors les longituds d'ona també seran la meitat:

$\lambda_A = 2L = 0,400\ \mathrm{m}$, $\lambda_B = L = 0,200\ \mathrm{m}$, $\lambda_C = \frac{2}{3}L = 0,133\ \mathrm{m}$ i $\lambda_D = \frac{1}{2}L = 0,100\ \mathrm{m}$.

\punts{0,6} Com $v = \lambda\cdot f \Rightarrow f = \frac{v}{\lambda}$, per tant les freqüències seran el doble:

$f_A = 240\ \mathrm{Hz}$, $f_B = 480\ \mathrm{Hz}$, $f_C = 720\ \mathrm{Hz}$ i $f_D = 960\ \mathrm{Hz}$.
\end{solapartat}
